% Part: sets-functions-relations
% Chapter: arithmetization
% Section: cuts
%
\documentclass[../../../include/open-logic-section]{subfiles}

\begin{document}

\olfileid[fa-IR]{sfr}{arith}{cuts}
\olsection{از $\Rat$ به $\Real$}

در اصل، خاصیت تمامیت نشان می‌دهد که هر نقطهٔ $\alpha$ از خط حقیقی، آن
خط را به‌طور کامل به دو نیمه تقسیم می‌کند: نیمه‌ای که $\alpha$ کوچک‌ترین
کران بالای آن است، و نیمه‌ای که $\alpha$ بزرگ‌ترین کران پایین آن است.
ددکیند پیشنهاد کرد که برای \emph{ساختن} اعداد حقیقی از اعداد گویا، اعداد
حقیقی را همان \emph{برش‌هایی} بدانیم که اعداد گویا را افراز می‌کنند. یعنی
$\sqrt{2}$ را با \emph{برشی} یکی می‌گیریم که اعداد گویای
$< \sqrt{2}$ را از اعداد گویای $> \sqrt{2}$ جدا می‌کند.

بیایید این مطلب را دقیق‌تر کنیم. اگر اعداد گویا را به دو نیمه ببُریم، تنها
با درنظرگرفتن نیمهٔ \emph{پایینی} می‌توانیم افرازی را که ساخته‌ایم به‌طور
یکتا مشخص کنیم. بنابراین، برای بیان دقیق، تعریف زیر را عرضه می‌کنیم:

\begin{defn}[برش]
یک \emph{برش} $\alpha$ هر پارهٔ آغازینِ ناتهی و حقیقی از اعداد گویاست
که بزرگ‌ترین عضو ندارد. یعنی $\alpha$ برش است اگر و تنها اگر:
\begin{enumerate}
	\item \emph{ناتهی و حقیقی}: $\emptyset \neq \alpha \subsetneq \Rat$
	\item \emph{آغازین}: برای همهٔ $p,q \in \Rat$: اگر $p < q \in \alpha$، آنگاه $p \in \alpha$
	\item \emph{بدون بیشینه}: برای هر $p \in \alpha$، عددی $q \in \alpha$ وجود دارد چنان‌که $p < q$
\end{enumerate}
سپس $\Real$ مجموعهٔ برش‌هاست.
\end{defn}

اکنون می‌توانیم بگوییم $\sqrt{2} = \Setabs{p \in \Rat}{p^2 < 2\text{
یا }p < 0}$. البته باید بررسی کنیم که این مجموعه واقعاً \emph{برش} است،
اما این کار را به \olref[check]{sec} واگذار می‌کنیم.

مانند پیش، پس از تعریف چند موجودیت باید توابع و روابط پایه روی آن‌ها را
تعریف کنیم. با موردی آسان آغاز می‌کنیم:
\begin{align*}
	\alpha \leq \beta \text{ اگر و تنها اگر }\alpha \subseteq \beta
\end{align*}
این تعریف ترتیب به ما اجازه می‌دهد نتیجهٔ اصلی، یعنی برخورداری مجموعهٔ
برش‌ها از خاصیت تمامیت، را \emph{بیان} کنیم. صورت کامل گزاره چنین است.
اگر $S$ مجموعه‌ای ناتهی از برش‌ها باشد و کران بالا داشته باشد، آنگاه $S$
کوچک‌ترین کران بالا دارد. با جزئیات بیشتر: برشی چون $\lambda$ وجود دارد
که کران بالای $S$ است، یعنی\ $(\forall \alpha \in S)\alpha \subseteq
\lambda$، و $\lambda$ کوچک‌ترین چنین برشی است، یعنی\ $(\forall \beta \in
\Real)((\forall \alpha \in S)\alpha \subseteq \beta\lif \lambda
\subseteq \beta)$. اکنون برهان نتیجه را می‌آوریم:

\begin{thm}\ollabel{realcompleteness}
مجموعهٔ برش‌ها خاصیت تمامیت دارد.
\end{thm}

\begin{proof}
فرض کنید $S$ مجموعه‌ای ناتهی از برش‌ها باشد که کران بالا دارد. بگذارید
$\lambda = \bigcup S$.
%\Setabs{p \in \Rat}{(\exists \alpha \in S)p \in \alpha}$$
نخست ادعا می‌کنیم $\lambda$ یک برش است:
\begin{enumerate}
\item چون $S$ کران بالا دارد، دست‌کم یک برش در $S$ است، پس
$\emptyset \neq \lambda$. چون $S$ مجموعه‌ای از برش‌هاست، $\lambda
\subseteq \Rat$. چون $S$ کران بالا دارد، عددی $p \in \Rat$ هست که در
هیچ برش $\alpha \in S$ حضور ندارد. پس $p\notin \lambda$، و در نتیجه
$\lambda \subsetneq \Rat$.
\item فرض کنید $p < q \in \lambda$. پس $\alpha \in S$ای وجود دارد که
$q \in \alpha$. چون $\alpha$ برش است، $p \in \alpha$. بنابراین
$p \in \lambda$.
\item فرض کنید $p \in \lambda$. پس $\alpha \in S$ای وجود دارد که
$p \in \alpha$. چون $\alpha$ برش است، $q \in \alpha$ای وجود دارد که
$p < q$. بنابراین $q \in \lambda$.
\end{enumerate}
این ادعا را ثابت می‌کند. افزون بر این، آشکارا $(\forall \alpha \in
S)\alpha \subseteq \bigcup S = \lambda$؛ یعنی\ $\lambda$ کران بالای
$S$ است. اکنون فرض کنید $\beta \in \mathbb{R}$ نیز کران بالا باشد، یعنی\ $(\forall
\alpha \in S)\alpha \subseteq \beta$. برای هر $p \in \Rat$،
اگر $p \in \lambda$، آنگاه $\alpha \in S$ای هست که $p \in \alpha$،
پس $p \in \beta$. با تعمیم، $\lambda \subseteq \beta$. بنابراین
$\lambda$ \emph{کوچک‌ترین} کران بالای $S$ است.
\end{proof}

پس مجموعه‌ای از موجودیت‌ها داریم که خاصیت تمامیت را برآورده می‌کنند. یک
راه بیان این مطلب آن است که در برش‌های ما هیچ «شکافی» وجود ندارد. (بنابراین،
گرفتن «برش‌های» بیشتر از اعداد حقیقی، به‌جای اعداد گویا، هیچ شیء تازهٔ
جالبی به دست نمی‌دهد.)

اکنون باید چند عمل را روی اعداد حقیقی تعریف کنیم. با نشاندن اعداد گویا در
اعداد حقیقی آغاز می‌کنیم: مقرر می‌داریم $p_\Real = \Setabs{q \in
\Rat}{q < p}$، برای هر $p \in \Rat$. سپس تعریف می‌کنیم:
\begin{align*}
	\alpha + \beta &= \Setabs{p + q}{p \in \alpha \land q \in \beta}\\
%	\alpha - \beta &\defis \Setabs{p - q}{p \in \alpha \land q \in \Rat \setminus \beta}\\
	\alpha \times \beta &=
	\Setabs{p \times q}{0 \leq p \in \alpha \land 0 \leq q \in \beta} \cup 0^\mathbb{R} & \text{اگر }\alpha, \beta \geq 0_\Real
%	\alpha \div \beta &\defis 
%	\Setabs{\nicefrac{p}{q}}{p \in \alpha \land q \in \Rat \setminus \beta}  & \text{if }\alpha \geq 0_\Real, \beta > 0_\Real
\end{align*}
برای رسیدگی به دیگر حالت‌های ضرب، نخست بگذارید: %that $0_\Real \times \alpha = 0_\Real = \alpha \times 0_\Real$, and add:
\begin{align*}
	-\alpha &= \Setabs{p - q}{p < 0 \land q \notin \alpha}
\end{align*}
و سپس مقرر می‌داریم:
\begin{align*}
	\alpha \times \beta &\defis
	\begin{cases}
		\mathord{-}\alpha \times \mathord{-}\beta &\text{اگر }\alpha < 0_\Real\text{ و }\beta < 0_\Real\\
		\mathord{-}(\mathord{-}\alpha \times \beta) &\text{اگر }\alpha < 0_\Real \text{ و }\beta > 0_\Real\\
		\mathord{-}(\alpha \times \mathord{-}\beta) &\text{اگر }\alpha > 0_\Real \text{ و }\beta < 0_\Real
	\end{cases}
\end{align*}
سپس باید بررسی کنیم که هر یک از این تعریف‌ها همواره برشی به دست می‌دهد.
و در پایان باید نمایشی آسان (اما طولانی) عرضه کنیم که برش‌هایی که بدین
شیوه تعریف شده‌اند دقیقاً چنان رفتار می‌کنند که باید. همهٔ این کارها را به
\olref[check]{sec} واگذار می‌کنیم.
\end{document}
